\section*{Capitulo 8. Limites, continuidad y asintotas}

\begin{shortanswer}{[GX-C08.S01-01]}
\[
D_g=(-\infty,3).
\]
\end{shortanswer}

\begin{shortanswer}{[GX-C08.S01-02]}
\[
D_h=\R\setminus\{-3,3\}.
\]
\end{shortanswer}

\begin{shortanswer}{C08.S01 practica autonoma}
\begin{enumerate}[label=\textbf{PX-C08.S01-\arabic*:}, leftmargin=*]
  \item \(\R\)
  \item \(\R\setminus\{5\}\)
  \item \([4,\infty)\)
  \item \((-\infty,-2)\cup[1,\infty)\)
  \item \((-\infty,-2)\cup(2,\infty)\)
  \item \([-3,-1)\cup(-1,3]\)
\end{enumerate}
\end{shortanswer}

\begin{shortanswer}{[GX-C08.S02-01]}
Raices: \(-2\) y \(1\). Valor maximo: \(5\).
\end{shortanswer}

\begin{shortanswer}{[GX-C08.S02-02]}
\[
f(2)=1,\qquad \lim_{x\to 2}f(x)=3.
\]
\end{shortanswer}

\begin{shortanswer}{C08.S02 practica autonoma}
\begin{enumerate}[label=\textbf{PX-C08.S02-\arabic*:}, leftmargin=*]
  \item \(f(0)=3\). Raices: \(-2\) y \(1\)
  \item Maximo absoluto \(4\). Minimo absoluto \(-1\)
  \item En \(x=0\)
  \item No existe
  \item \(f(3)=2\) y \(\lim_{x\to 3}f(x)=-1\)
  \item La funcion tiende a una asintota horizontal \(y=4\)
\end{enumerate}
\end{shortanswer}

\begin{shortanswer}{[GX-C08.S03-01]}
\[
1.
\]
\end{shortanswer}

\begin{shortanswer}{[GX-C08.S03-02]}
\[
\lim_{x\to 0^-}\frac{1}{x}=-\infty,\qquad
\lim_{x\to 0^+}\frac{1}{x}=+\infty.
\]
\end{shortanswer}

\begin{shortanswer}{C08.S03 practica autonoma}
\begin{enumerate}[label=\textbf{PX-C08.S03-\arabic*:}, leftmargin=*]
  \item \(7\)
  \item \(2\)
  \item No existe
  \item \(3\)
  \item \(-\infty\) por la izquierda y \(+\infty\) por la derecha
  \item No
\end{enumerate}
\end{shortanswer}

\begin{shortanswer}{[GX-C08.S04-01]}
\[
5.
\]
\end{shortanswer}

\begin{shortanswer}{[GX-C08.S04-02]}
\[
0.
\]
\end{shortanswer}

\begin{shortanswer}{C08.S04 practica autonoma}
\begin{enumerate}[label=\textbf{PX-C08.S04-\arabic*:}, leftmargin=*]
  \item \(3\)
  \item \(0\)
  \item \(-\infty\)
  \item \(2\)
  \item \(0\)
  \item \(\frac{1}{3}\)
\end{enumerate}
\end{shortanswer}

\begin{shortanswer}{[GX-C08.S05-01]}
\[
6.
\]
\end{shortanswer}

\begin{shortanswer}{[GX-C08.S05-02]}
\[
\frac{1}{4}.
\]
\end{shortanswer}

\begin{shortanswer}{C08.S05 practica autonoma}
\begin{enumerate}[label=\textbf{PX-C08.S05-\arabic*:}, leftmargin=*]
  \item \(2\)
  \item \(-1\)
  \item \(\frac{1}{2}\)
  \item \(\frac{1}{4}\)
  \item \(-1\)
  \item \(\frac{1}{2}\)
\end{enumerate}
\end{shortanswer}

\begin{shortanswer}{[GX-C08.S06-01]}
Vertical: \(x=3\). Horizontal: \(y=2\).
\end{shortanswer}

\begin{shortanswer}{[GX-C08.S06-02]}
Vertical: \(x=-2\). Horizontal: \(y=0\).
\end{shortanswer}

\begin{shortanswer}{C08.S06 practica autonoma}
\begin{enumerate}[label=\textbf{PX-C08.S06-\arabic*:}, leftmargin=*]
  \item Vertical \(x=-4\). Horizontal \(y=3\)
  \item Horizontal \(y=0\)
  \item No tiene asintota vertical; la grafica coincide con \(y=x+2\) salvo en \(x=2\)
  \item No tiene asintota vertical; la grafica coincide con \(y=x+3\) salvo en \(x=0\)
  \item Vertical \(x=-1\). Oblicua \(y=x-1\)
  \item Verticales \(x=-3\) y \(x=3\). Horizontal \(y=0\)
\end{enumerate}
\end{shortanswer}

\begin{shortanswer}{[GX-C08.S07-01]}
No es continua en \(x=0\). Tiene una discontinuidad de salto.
\end{shortanswer}

\begin{shortanswer}{[GX-C08.S07-02]}
Discontinuidad infinita.
\end{shortanswer}

\begin{shortanswer}{C08.S07 practica autonoma}
\begin{enumerate}[label=\textbf{PX-C08.S07-\arabic*:}, leftmargin=*]
  \item Si, es continua
  \item Removible
  \item Infinita
  \item De salto
  \item Si, es continua
  \item Si
\end{enumerate}
\end{shortanswer}

\begin{shortanswer}{[GX-C08.S08-01]}
No es continua en \(x=1\). Hay un salto.
\end{shortanswer}

\begin{shortanswer}{[GX-C08.S08-02]}
\[
g(2)=0
\]
y es continua en \(x=2\).
\end{shortanswer}

\begin{shortanswer}{C08.S08 practica autonoma}
\begin{enumerate}[label=\textbf{PX-C08.S08-\arabic*:}, leftmargin=*]
  \item Si, es continua
  \item Si, es continua
  \item \(4\)
  \item Tiene un pico, no un salto
  \item Si, es continua
  \item Semirrecta horizontal \(y=1\) a la izquierda y recta \(y=x+1\) a la derecha, unidas en \((0,1)\)
\end{enumerate}
\end{shortanswer}

\begin{shortanswer}{[GX-C08.S09-01]}
\[
k=2.
\]
\end{shortanswer}

\begin{shortanswer}{[GX-C08.S09-02]}
\[
m=1.
\]
\end{shortanswer}

\begin{shortanswer}{C08.S09 practica autonoma}
\begin{enumerate}[label=\textbf{PX-C08.S09-\arabic*:}, leftmargin=*]
  \item \(a=3\)
  \item \(b=4\)
  \item \(c=4\)
  \item \(d=2\)
  \item \(k=6\)
  \item \(m=1\)
\end{enumerate}
\end{shortanswer}

\begin{shortanswer}{[GX-C08.S10-01]}
\[
80.
\]
\end{shortanswer}

\begin{shortanswer}{[GX-C08.S10-02]}
\[
5.
\]
\end{shortanswer}

\begin{shortanswer}{C08.S10 practica autonoma}
\begin{enumerate}[label=\textbf{PX-C08.S10-\arabic*:}, leftmargin=*]
  \item \(30\)
  \item \(3\)
  \item Inicial \(85\). Limite \(100\)
  \item \(20\)
  \item \(0\)
  \item \(6\) y \(4\)
\end{enumerate}
\end{shortanswer}

\begin{shortanswer}{[GX-C08.S11-01]}
\[
4.
\]
\end{shortanswer}

\begin{shortanswer}{[GX-C08.S11-02]}
\[
a=2.
\]
\end{shortanswer}

\begin{shortanswer}{C08.S11 practica autonoma}
\begin{enumerate}[label=\textbf{PX-C08.S11-\arabic*:}, leftmargin=*]
  \item \(1\)
  \item No existe como numero finito; diverge
  \item \(k=2\)
  \item Vertical \(x=-2\). Oblicua \(y=x-2\)
  \item \(1\)
  \item No
\end{enumerate}
\end{shortanswer}
